\section{Methods: Local Ancestry Inference via Multi-State HMM}
\label{sec:methods_ancestry}

\subsection{Model Formulation}

We extend the HMM framework to an $N$-state model where each hidden state corresponds to one of $N$ reference populations. Given pairwise similarity between a query haplotype and reference panel haplotypes, the model infers the most likely ancestral origin at each genomic window. This performs a dimensional reduction relative to the Li \& Stephens copying model~\cite{li2003stephens} underlying RFMix~\cite{maples2013rfmix}: instead of tracking which individual reference haplotype is copied at each SNP (requiring $>$500 haplotypes), we ask which \emph{population} shows highest aggregate identity per window. This collapses the state space from $m$ haplotypes to $k$ populations; empirically, $\sim$50 haplotypes suffice versus $>$1{,}400 for haplotype-matching methods.

\paragraph{Emission model.}
For each population $k$, per-haplotype similarities are aggregated into a single score $f_k(o_t)$ using one of four functions: max (default), mean, median, or top-$j$ averaging. The emission probability uses a softmax parameterized by temperature $\tau$:
\begin{equation}
P(o_t \mid s_t = k) = \frac{\exp(f_k(o_t) / \tau)}{\sum_{j : f_j > 0} \exp(f_j(o_t) / \tau)}
\end{equation}
The softmax is the Bayes-optimal posterior under Gaussian discriminant analysis with shared variance, and satisfies Luce's choice axiom. The temperature $\tau$ is estimated adaptively from the median per-window spread in population scores, with theoretical scaling $\tau_{\text{opt}} \propto 1/\sqrt{\ln n}$ by extreme value theory.

\paragraph{Pairwise contrast emissions.}
\label{sec:pairwise_contrast}
As an alternative to joint softmax, each population pair $(i, j)$ is evaluated independently via a Bradley-Terry comparison:
\begin{equation}
P_{ij}(k = i \mid o_t) = \frac{\exp(f_i / \tau_{ij})}{\exp(f_i / \tau_{ij}) + \exp(f_j / \tau_{ij})}
\end{equation}
with per-pair adaptive temperature $\hat{\tau}_{ij} = \text{median}_t |f_i(o_t) - f_j(o_t)|$. The per-population emission is obtained by aggregating across all pairwise comparisons: $P(o_t \mid k) \propto \prod_{j \neq k} P_{kj}(k \mid o_t)$. This resolves the problem that a single temperature cannot simultaneously be sharp for strong contrasts (AFR--EUR, $F_{ST} \approx 0.1$) and gentle for weak ones (EUR--AMR, $F_{ST} \approx 0.02$). Empirically, pairwise contrast improves simulated ancestry from 95.1\% (joint softmax) to 97.95\%.

\paragraph{Transitions.}
The transition matrix is uniform across populations, parameterized by a single switch probability $p_{\text{switch}}$. The initial state distribution is uniform. When a genetic map is provided, $p_{\text{switch}}$ is modulated by local recombination rate via the Haldane map function. The switch probability is estimated in a two-pass approach: an initial Viterbi pass with a broad prior, followed by regularization toward $p_{\text{switch}} = 0.001$.

\subsection{Automatic Parameter Configuration}
\label{sec:auto_configure}

Optimal parameters differ substantially between datasets: simulated data with strong population separation requires $w = 0.7$, $ec = 1$, while HPRC real data requires $w = 0.12$, $ec = 15$, with a 9.8\,pp cross-application penalty. We introduce automatic configuration based on $D_{\min}$, the minimum Cohen's $d$ across all reference population pairs:
\begin{equation}
D_{\min} = \min_{i < j} \frac{|\bar{f}_i - \bar{f}_j|}{\sqrt{(\hat{\sigma}_i^2 + \hat{\sigma}_j^2)/2}}
\end{equation}
The pairwise weight scales with the coefficient of variation of pairwise Cohen's $d$ values, while the emission context scales inversely with $D_{\min}$: $ec^* = \text{round}(0.1/D_{\min})$, clamped to $[1, 15]$. This formula recovers 3.3\,pp on simulated data and eliminates manual parameter tuning.

\subsection{Inference}

The most likely ancestry sequence is computed via Viterbi decoding; posterior ancestry probabilities via forward-backward. As an alternative, posterior decoding ($\hat{a}_t = \arg\max_k \gamma_t(k)$) makes independent per-window decisions, which is advantageous for detecting short minority ancestry tracts where Viterbi's path constraint suppresses isolated switches. The switch probability is refined via Baum-Welch training over all sample sequences simultaneously. Genome-wide admixture proportions, per-population tract lengths, ancestry LOD scores, cross-validation, and output format details are in Supplementary Section~L.4.
